Fix a nonzero zero-sum contrast \(c\) with \(|S_c|\geq3\). The established first-order upper bound and embedded two-arm converse imply
\[
 0\leq d_n(c)\leq C_0(c)\{n^{-1}-\rho_{n,2}\}
 =C_0(c)C_A n^{-4/3}+o(n^{-4/3}).
\]
The support-two case is excluded from this question: when \(|S_c|=2\), equality holds throughout the preceding display, so \(a_n=n^{4/3}\) and \(C_{\mathrm{so}}(c)=C_0(c)C_A\) are already supplied by the published two-arm theorem and its exact embedding. In the genuinely multivariate regime \(|S_c|\geq3\), use the active-face boundary-layer handle to determine whether a regularly varying normalizer \(a_n\to\infty\) exists for which \(a_nd_n(c)\) converges to a finite positive limit. In the existence branch, characterize the limiting value problem and prove existence and uniqueness of its admissible solution; in the alternative branch, prove that no regularly varying normalization has a finite positive limit and identify the correct asymptotic behavior. Decide from proof whether the limit is a multidimensional free-boundary Hamilton--Jacobi--Bellman problem, an elliptic cone spectral problem, or a separable collection of scalar layers. For \(c=c^\dagger\), any conclusion must agree with certified exact-rational finite-\(n\) bounds, which are evidence only and do not by themselves prove \(n^{4/3}\) convergence, positivity of a new constant, separability, or a multidimensional continuum operator.